\section{Common isotropic three-dimensional subspaces}\label{sec:subspaces}

The linear algebra in this section is over an arbitrary field of
characteristic zero. Nodes are distinct, $\gamma_i$ is defined by
\eqref{eq:barycentric}, and $W$ is the quadratic evaluation space.
A subspace is \emph{totally isotropic} for a symmetric bilinear form
if that form vanishes on every pair of its vectors.

\begin{lemma}[Seven nodes]\label{lem:seven}
For seven nodes, every three-dimensional subspace totally isotropic
for $B_0$ and $B_1$ is
\begin{equation}\label{eq:sign-space}
 U=E W,\qquad E=\diag(\eps_1,\ldots,\eps_7),\quad \eps_i\in\{1,-1\}.
\end{equation}
There are exactly $2^6$ such spaces.
\end{lemma}
\begin{proof}
Let $T=\diag(t_i)$. Since $B_1(a,b)=B_0(Ta,b)$, we have
$TU\subset U^{\perp_{B_0}}$. The latter has dimension four and
contains $U$. Thus
$K=\{u\in U:Tu\in U\}$ has dimension at least two. The linear map
$K\to U/K$, $u\mapsto Tu\bmod K$, has a nonzero kernel. Choose
$u\ne0$ in it; then $u,Tu,T^2u\in U$.

These vectors are independent. Otherwise $f(T)u=0$ for some nonzero
polynomial of degree at most two, so $u$ is supported on at most two
nodes. But $B_0(u,u)=B_1(u,u)=0$ is an invertible Vandermonde
system on those nodes for the quantities $\gamma_i u_i^2$. It
would force $u=0$.

The three vectors therefore form a basis of $U$. Applying both
forms to their pairs yields
\[
 \sum_i\gamma_i t_i^r u_i^2=0\qquad(0\le r\le5).
\]
The $6\times7$ Vandermonde matrix has a one-dimensional kernel,
spanned by $(\gamma_i)_i$ by \eqref{eq:moments}. Hence
$u_i^2=c$ for every $i$. Since $u\ne0$, $c\ne0$. Dividing by
$u_1$ gives $u_i/u_1=\eps_i\in\{1,-1\}$ and proves
\eqref{eq:sign-space}. Conversely every such space is isotropic
by \eqref{eq:moments}.

For the count, suppose a relative sign matrix preserves $W$.
Its sign vector is the evaluation of a polynomial $q$ of degree at
most two, and its product with $(t_i^2)_i$ is the evaluation of
another such polynomial $r$. The polynomial $T^2q(T)-r(T)$ has
degree at most four and seven roots, so is zero. Thus $q$ is
constant, and all relative signs coincide. Only a common sign
is redundant, giving $2^6$ spaces.
\end{proof}

\begin{lemma}[Eight nodes]\label{lem:eight}
For eight nodes, a three-dimensional subspace totally isotropic
for $B_0,B_1,B_2$ is $EW$ for a sign matrix $E$.
\end{lemma}
\begin{proof}
The map $T$ is injective on $U$. Indeed a vector in its kernel
would be supported at the unique possible zero node, and a
nonzero such vector has $B_0(u,u)\ne0$.
The subspace $U+TU$ is totally isotropic for $B_0$: the
$U,U$ pair uses $B_0$, the mixed pair uses $B_1$, and the
$TU,TU$ pair uses $B_2$. Nondegeneracy of $B_0$ in dimension
eight gives $\dim(U+TU)\le4$. Since $\dim TU=3$, the space
$K=\{u\in U:Tu\in U\}$ again has dimension at least two.
The preceding kernel argument supplies $u,Tu,T^2u\in U$;
their independence follows as in Lemma~\ref{lem:seven}.
The three forms on these vectors give
\[
 \sum_i\gamma_i t_i^r u_i^2=0\qquad(0\le r\le6).
\]
The $7\times8$ Vandermonde kernel is spanned by $(\gamma_i)_i$,
so all $u_i^2$ are equal and nonzero. Normalization gives $U=EW$.
The sign-uniqueness argument in Lemma~\ref{lem:seven} also
applies to eight nodes; there are $2^7$ such spaces.
\end{proof}

\begin{proposition}[Restriction and gluing]\label{prop:gluing}
Let $n\ge7$. Suppose $U\subset F^n$ has dimension three and
satisfies the product property \eqref{eq:product-property}.
Then $U=EW$ for a sign matrix $E$.
\end{proposition}
\begin{proof}
For $n=7$, the product property implies isotropy for $B_0,B_1$
by interpolation, so use Lemma~\ref{lem:seven}. Suppose $n\ge8$.
Choose eight coordinate positions such that restriction of $U$
has rank three: choose three independent coordinate rows of a
basis matrix and add five positions. The restricted product vectors
are quartic evaluations. Interpolation on these eight nodes makes
the restricted space isotropic for its own $B_0,B_1,B_2$.
Lemma~\ref{lem:eight} makes it a signed quadratic evaluation
space. In particular, any three of its coordinate rows are
independent.

For any additional node, combine it with any seven of the chosen
eight. The new restriction has rank three and is again a signed
quadratic space. On the common seven nodes the two spaces agree:
both are the image of $U$ there, and restriction has rank three.
The uniqueness of signs in Lemma~\ref{lem:seven} makes their
signs agree up to one common reversal. After this reversal, a
fixed vector of $U$ is represented on both sets by quadratic
polynomials agreeing at seven nodes; they are the same polynomial.
Thus the new sign extends the original signs consistently for
every vector in $U$. Repeating for each remaining node proves
$U=EW$ globally.
\end{proof}

\begin{lemma}[The two sign groups]\label{lem:sign-groups}
Suppose $n\ge7$, $\dim V=3$, $V\cap W=0$, and
$V\subset W+EW$. Then each sign occurs at least three times.
Moreover, restriction of every vector of $V$ to either sign group
is the evaluation of a quadratic polynomial on that group.
\end{lemma}
\begin{proof}
Let $m$ be the size of the smaller sign group $I$ and let $J$ be
the other group. If $m\le2$, then $|J|\ge5$.
For a vector in $W\cap EW$, write its values as $r(t_i)$ and
$\eps_iq(t_i)$ and reverse all signs so that $\eps_i=1$ on $J$.
On $J$ the quadratics $r,q$ agree, hence are identical. On $I$
the polynomial $q$ must vanish. Therefore
$\dim(W\cap EW)=3-m$ and
\[
 \dim((W+EW)/W)=m<3.
\]
This contradicts the injective image of $V$ in the quotient.
Both groups consequently have size at least three.
Finally, an element of $W+EW$ has the form
$(p(t_i)+\eps_iq(t_i))_i$. Its restrictions are the quadratics
$p+q$ and $p-q$, respectively.
\end{proof}

All these classifications are field arguments, rather than
enumerations of rational examples. As arithmetic checks,
\ancfile{structure.py} verifies the $64$ seven-node sign spaces
and their projection dimensions, and \ancfile{eight_factory.py}
checks the $128$ eight-node sign spaces and $\dim(U+TU)=4$
on distinct rational nodes. These checks supplement the proofs
above.
